\subsection{Special Forms: Difference of Squares} 
\begin{Example}
Factor $x^2-9$ using the AC Method.
\begin{itemize*}
\item To apply the AC Method, we need to add the term \makeblanksmall.
\end{itemize*}
\begin{lpic}{}{4}
\end{lpic}
\end{Example}
\noi We can generalize this to a form called \term{difference of squares} by multiplying:
\begin{lpic}{}{1}
\regtextleft{0.5}{-1}{$(A-B)(A+B)={}$};
\end{lpic}
\begin{fact}[Recognizing a Difference of Squares]
A binomial is a \term{difference of squares} when:
\begin{itemize*}
\item Each term is a perfect square monomial (ignoring signs).
\item One is positive, and the other negative.
\end{itemize*}
\end{fact}
\begin{Example}
For each polynomial below:
\begin{enumerate*}\setabc
\item Determine whether it is a difference of squares.
\item If it is, factor it as such.
\end{enumerate*}
\begin{lpicsplit}{}{8}
\foreach \x/\exp in {0/9x^2-16,1/x^4y^2-36z^6} {
	\regtextleft{0.5+\x*\value{fcols}+\x}{-1}{$\exp$};
	\foreach \y/\letter in {-2/a,-3/b}
		\regtextright{1+\x*\value{fcols}+\x}{\y}{\letter.};
}
\end{lpicsplit}
\begin{lpicsplit}{}{3}
\foreach \x/\exp in {1/x^3-y^2,0/z^4+81} {
	\regtextleft{0.5+\x*\value{fcols}+\x}{-1}{$\exp$};
	\foreach \y/\letter in {-2/a,-3/b}
		\regtextright{1+\x*\value{fcols}+\x}{\y}{\letter.};
}
\end{lpicsplit}
\end{Example}